\section{Conclusion and Limitations}
The layer that ranks a set of tokens highly need not be the layer at which deleting them does the most damage. For both late-selected sets, removing the same tokens at L8 hurt more than removing them at L14. A pruning evaluation should therefore report two depths: where the ranking was obtained and where the intervention was applied. When the goal was instead to retain useful tokens, L14/128 outperformed two early schedules at matched analytical LM cost on the ScanQA validation split used for selection. Whether that advantage holds on unseen scenes remains open: our transfer tests compare against the unpruned models, not cost-matched early schedules.

The study covers two Llama-family 7B backbones. At the default retention budget, quality declines on some tasks, and on one H100 the application-time savings are much smaller than the analytical LM savings. We did not measure application time against official 3DZip. SQA3D test was previously inspected, and historical training exposure remains incompletely verified. The supplement details limits on semantic interpretation, mixed-batch timing intervals, and reproduction of the historical quality run.\label{maintext:end}
